%---------------------------------------------------------------------
\subsection{Discussion}

\subsubsection{Evidence Synthesis}

The intervention trials, one meta-analysis, and one injury survey support five conclusions with the confidence levels noted:
\begin{enumerate}[noitemsep, topsep=2pt]
  \item Hangboard supplementation consistently improves finger-specific MFS and endurance in intermediate-to-advanced male climbers within 4--10 weeks (meta-analytic SDM 0.41--1.23 \cite{SO23}); evidence for female and beginner populations is very limited \cite{LAS24}.
  \item Intensity determines adaptation type: $\geq80\%$ MFS maximises MFS; 60--80\% MFS maximises stamina and endurance; 80\% MFS improves all three \cite{DO22}, making it the pragmatic default for athletes who cannot dedicate separate training blocks.
  \item Combining max-hang and repeater protocols within one mesocycle underperforms either in isolation on endurance \cite{O19}, consistent with concurrent-training interference, though this mechanism is untested in climbing.
  \item Progressive load on a standard edge outperforms progressive edge reduction for grip strength (one RCT, small ES 0.36) \cite{MSAE21}; community preference for minimal-edge work is not supported by direct evidence.
  \item Daily low-intensity hangs (Abrahangs, ${\sim}40\%$ MFS) did not significantly outperform climbing-only controls in the only available data ($+2.5\%$, $p = 0.12$), while twice-weekly max hangs did ($+3.2\%$, $p < 0.001$); combining both was additive ($+5.8\%$) \cite{GKAO24}. The design is retrospective observational; findings are hypothesis-generating only. Fingerboard training is associated with lower injury risk in experienced climbers ($\geq6$ yr) but higher risk in rapidly advancing less-experienced climbers \cite{O23}.
\end{enumerate}

\subsubsection{Open Questions}

\begin{itemize}[noitemsep, topsep=2pt]
  \item \textbf{Grade transfer}: No study demonstrates statistically significant redpoint-grade improvement; the Stien meta found no effect on climbing-specific on-wall tests \cite{SO23}. Dynamometric gains $\neq$ climbing performance gains.
  \item \textbf{Beginner and female populations}: Sjöman \cite{O23} suggests elevated injury risk in rapidly advancing less-experienced climbers; Langer \cite{LAS24} provides the only female RCT data (anecdotal effects). Both gaps are critical.
  \item \textbf{Long-term tendon adaptation}: No trial exceeds 10 weeks; collagen remodelling operates on a timescale of months. No study provides tendon imaging data.
  \item \textbf{Optimal periodisation}: Sequencing of max-hang and repeater blocks, and integration of daily low-intensity loading, across a training year is entirely untested.
  \item \textbf{Abrahangs RCT}: The observational Gilmore cohort \cite{GKAO24} is hypothesis-generating; a direct RCT with controlled loads and compliance is needed before recommending daily loading to all populations.
\end{itemize}

\subsubsection{Evidence-Weighted Protocol}

For intermediate-to-advanced climbers with $\geq1$--2 years of consistent climbing, no active finger injury, and a male-cohort evidence base (generalisation to females and beginners is unsupported). All gains documented here are on dynamometric surrogates; grade improvement is unproven.

\begin{itemize}[noitemsep, topsep=2pt]
  \item \textbf{Edge}: 18--20~mm half-crimp; increase intensity by adding weight before reducing depth \textnormal{[RCT: \cite{MSAE21}, ES 0.36]}.
  \item \textbf{Strength phase (off-season)}: Max hangs at 85--95\% MFS, 5--10~s, 3--6 sets, 2--3~min rest, 2$\times$/week with $\geq48$~h recovery \textnormal{[RCT: \cite{O19, MSAE21}; meta: \cite{SO23}]}.
  \item \textbf{Stamina phase (pre-season)}: Repeaters at 60--80\% MFS, 7~s on / 3~s off $\times$ 6 reps per set, 2$\times$/week \textnormal{[RCT: \cite{HO22, DO22}]}.
  \item \textbf{Daily loading (experimental, observational only)}: Abrahangs at ${\sim}40\%$ MFS did not significantly outperform climbing-only controls as a standalone protocol ($p = 0.12$); adding them to max hangs was additive ($+5.8\%$) in a retrospective cohort \textnormal{[obs.: \cite{GKAO24}]}. No RCT has tested this; treat as experimental. Load self-assessment is unreliable; monitor for cumulative overuse.
  \item \textbf{Deload}: Every 4--8~weeks, reduce volume or intensity ${\sim}40\%$ for one week \textnormal{[practitioner consensus; untested in RCTs]}.
  \item \textbf{Injury monitoring}: Cease or reduce loading at first sign of SPIIF; less-experienced climbers ($<6$ yr) advancing rapidly face elevated injury risk with RFT \textnormal{[survey: \cite{O23}]}.
\end{itemize}
